\documentclass[10pt,journal]{IEEEtran}

\usepackage{booktabs}
\usepackage{graphicx}
\usepackage{amsmath}
\usepackage{siunitx}
\usepackage{xcolor}
\usepackage{url}
\usepackage{array}

\title{Scheduler-Derived Budget Interfaces for Dynamic Photonic Transformer Accelerators}
\author{Anonymous Authors}

\begin{document}
\maketitle

\begin{abstract}
Dynamic photonic tensor cores can accelerate Transformer attention and
feed-forward kernels, but their system efficiency is constrained by conversion,
memory movement, operand broadcast, and a missing control path between the
software scheduler and accelerator-internal quality knobs. This paper presents
SUDS, a scheduler-derived budget interface for Dynamic Photonic Transformer
accelerators. SUDS exposes timing slack as KEEP, DEGRADE, and PRUNE budgets for
Lightening-style dynamic photonic tensor-core tiles, then composes those
budgets with local L1 or signal-amplitude selectors for ADC precision,
degradation, and pruning decisions. The architecture model maps BERT-base GLUE
encoder kernels and MobileViT-S Transformer blocks to four-tile, two-core
DPTC-style optical arrays and accounts for optical compute, E/O and O/E
conversion, memory movement, optical-link cost, sideband-control overhead, and
sensitivity settings tied to ADC macro, RTL, and PHY boundary artifacts. Across
two Transformer workloads, the selected SUDS operating point keeps governed MPS
accuracy rows linked to architecture-modeled PPA and preserves at least
20.99\% pessimistic EDP improvement versus the same-scope Lightening-style
DPTC reference. The evidence is architecture simulation plus governed MPS
accuracy validation and calibration artifacts. It is not a fabrication, layout,
foundry, circuit-signoff, optical-device-signoff, or bench-energy claim.
\end{abstract}

\section{Introduction}

Transformer accelerators based on dynamic photonic tensor cores can attack the
matrix multiplications that dominate attention and feed-forward layers. A
Lightening-style DPTC fabric supplies a useful comparison point because it
supports dynamic full-range matrix multiplication rather than only
weight-static optical tensors. HyAtten-style systems sharpen a different
observation: conversion is often the limiting term, so low-resolution
converters and local signal selectors matter. These mechanisms still leave a
control-interface gap. The scheduler already knows when a kernel has timing
slack, but converter precision, degradation, and pruning choices are usually
selected by local model or signal proxies alone.

SUDS turns that timing information into a resource-control budget. The
scheduler exports per-layer and per-column slack to the accelerator control
plane. The photonic tile controller maps slack to KEEP, DEGRADE, and PRUNE
budgets. Local selectors then choose the exact columns inside those budgets.
This framing is important: SUDS is not a claim that slack alone is a semantic
importance oracle. It is a budget interface that composes scheduler pressure
with local L1, signal-amplitude, or overflow-style selectors.

The prior protected manuscript in this repository targets a methodology route.
This TETC source is a separate architecture-first route and does not overwrite
that fallback. The contribution stack is restricted to architecture modeling,
measured MPS accuracy linkage, and calibration evidence.

\section{Contributions}

This paper makes five architecture-facing contributions:

\begin{enumerate}
  \item A Transformer-to-DPTC mapping for BERT-base encoder kernels and
  MobileViT-S Transformer blocks, including conversion, memory, optical-link,
  and sideband-control accounting.
  \item A scheduler-derived SUDS budget interface that converts timing slack
  into explicit KEEP, DEGRADE, and PRUNE budgets for ADC precision,
  degradation, and pruning.
  \item A deterministic architecture simulator and design-space sweep over
  tile dimension, tile count, cores per tile, sideband grouping, and
  ADC-sharing mode.
  \item A matched comparison matrix covering uniform 8-bit, uniform 4-bit,
  random, L1, slack-only, SUDS-only, SUDS+L1, SUDS+signal, HyAtten-style,
  TeMPO-style, ASTRA-style, and Lightening-style DPTC rows.
  \item A calibration table tying ADC macro, RTL sideband, and PHY boundary
  artifacts to simulator parameters while keeping their evidence labels
  bounded.
\end{enumerate}

\section{Architecture Model}

\subsection{Transformer Kernel Mapping}

The simulator represents each Transformer operation as a tiled matrix
multiplication on a DPTC array. BERT-base uses the 12-layer encoder at sequence
length 128 and includes QKV projection, QK score, attention-value, output
projection, and two FFN kernels per layer. MobileViT-S uses the repository's
Transformer-block operation manifest and promotes the linear attention and FFN
kernels under \texttt{global\_rep}. Each GEMM is lowered into DPTC tile work,
with latency and energy accumulated over optical compute, DAC/MZM,
detector/TIA, ADC, laser, SRAM movement, optical-link movement, and
sideband-control groups.

\subsection{Selected Operating Point}

The selected operating point intentionally stays close to the Lightening and
HyAtten comparison surface. It uses 32 by 32 DPTC tiles, four tiles, two cores
per tile, 32-column sideband groups, and temporal accumulation before ADC. The
design-space artifact sweeps 3888 configurations and identifies 511 Pareto
rows; Table~\ref{tab:operating-point} records the promoted setting.

\begin{table}[t]
\centering
\caption{Selected architecture operating point. The choice keeps the main
claim on a Lightening-style DPTC fabric while exposing enough design-space
evidence for architecture review.}
\label{tab:operating-point}
\footnotesize
\begin{tabular}{@{}p{0.25\columnwidth}p{0.20\columnwidth}p{0.43\columnwidth}@{}}
\toprule
Knob & Selected value & Reason \\
\midrule
Tile dimension & \(32 \times 32\) & Matched DPTC comparison surface \\
Tiles & 4 & LT-B-style reference point \\
Cores per tile & 2 & Per-tile accumulation point \\
Sideband group & 32 columns & RTL control calibration anchor \\
ADC sharing & Temporal accumulation & Output-stationary DPTC readout \\
\bottomrule
\end{tabular}
\end{table}

\subsection{SUDS Control Layer}

For each scheduled kernel, the simulator records issue time, tile latency,
downstream deadline, and normalized slack. SUDS maps slack to a tier budget:
low slack keeps high-resolution conversion, medium slack degrades to
low-resolution conversion, and high slack permits pruning or stronger
degradation. SUDS+L1 and SUDS+signal keep the scheduler-derived tier counts but
use local selectors to choose exact columns. SUDS-only is retained as an
ablation that isolates the budget interface from the local selector.

\section{Evaluation Surface}

\subsection{Workloads and Evidence Labels}

The promoted architecture rows cover two Transformer workloads. BERT-base GLUE
accuracy rows are measured on PyTorch MPS and linked to the BERT DPTC schedule
through
\path{experiments/results/report_data/suds_glue_architecture_linkage_20260513_tetc_pivot.csv}.
MobileViT-S accuracy rows come from the governed MPS ImageNet validation
matrix. Energy, latency, area, optical-link cost, memory movement, and control
overhead remain architecture-modeled in both cases.

\begin{table*}[t]
\centering
\caption{Nominal selected-point PPA summary. Ratios are normalized to the
same-scope Lightening-style DPTC reference for each workload. Boundary rows are
shown for reviewer context, not as the selected SUDS fabric.}
\label{tab:ppa-summary}
\footnotesize
\begin{tabular}{llrrrr}
\toprule
Workload & Condition & Energy ratio & EDP ratio & Latency (ns) & Accuracy delta \\
\midrule
BERT-base GLUE & Lightening DPTC & 1.000 & 1.000 & 8526.22 & 0.00 pp \\
BERT-base GLUE & Uniform 4-bit & 0.976 & 0.976 & 8526.18 & n/a \\
BERT-base GLUE & L1 selector & 0.708 & 0.665 & 8015.11 & 0.00 pp \\
BERT-base GLUE & Slack-only & 0.708 & 0.666 & 8019.26 & 0.00 pp \\
BERT-base GLUE & SUDS+L1 & 0.785 & 0.754 & 8187.26 & 0.00 pp \\
BERT-base GLUE & SUDS+signal & 0.785 & 0.754 & 8187.26 & 0.00 pp \\
BERT-base GLUE & HyAtten-style boundary & 0.842 & 0.817 & 8274.59 & n/a \\
BERT-base GLUE & TeMPO boundary & 0.699 & 0.817 & 9961.96 & n/a \\
BERT-base GLUE & ASTRA boundary & 0.311 & 0.408 & 11198.49 & n/a \\
\midrule
MobileViT-S blocks & Lightening DPTC & 1.000 & 1.000 & 1193.95 & 0.00 pp \\
MobileViT-S blocks & Uniform 4-bit & 0.902 & 0.902 & 1193.92 & n/a \\
MobileViT-S blocks & L1 selector & 0.709 & 0.666 & 1122.33 & -3.47 pp \\
MobileViT-S blocks & Slack-only & 0.712 & 0.671 & 1125.17 & -2.66 pp \\
MobileViT-S blocks & SUDS+L1 & 0.775 & 0.746 & 1149.49 & -1.78 pp \\
MobileViT-S blocks & SUDS+signal & 0.775 & 0.746 & 1149.49 & -1.52 pp \\
MobileViT-S blocks & HyAtten-style boundary & 0.800 & 0.778 & 1160.29 & n/a \\
MobileViT-S blocks & TeMPO boundary & 0.720 & 0.842 & 1395.63 & n/a \\
MobileViT-S blocks & ASTRA boundary & 0.372 & 0.488 & 1568.86 & n/a \\
\bottomrule
\end{tabular}
\end{table*}

Table~\ref{tab:ppa-summary} shows the claim surface. SUDS+L1 and SUDS+signal
are promoted as scheduler-budgeted compositions. L1-only, slack-only,
signal-only, HyAtten-style, TeMPO-style, and ASTRA-style rows are retained so
reviewers can see when a local selector or alternate fabric is stronger. Those
rows bound the claim; they do not relabel SUDS as a universal winner.

\subsection{Pessimistic Gate}

The gate stresses ADC scaling, memory movement, optical-link penalties,
control overhead, and reduced compute relief from pruning. Under that setting,
the selected SUDS point preserves 18.58\% energy and 20.99\% EDP improvement
for BERT-base GLUE, and 21.84\% energy and 23.99\% EDP improvement for
MobileViT-S Transformer blocks, versus the same-scope Lightening-style
reference. The gate report stores this as the minimum required pessimistic EDP
margin.

\section{Baselines and Ablations}

\begin{table}[t]
\centering
\caption{Baseline and ablation contract. Main-text claims are made only where
the evidence status and claim boundary match the selected DPTC fabric.}
\label{tab:baseline-contract}
\footnotesize
\begin{tabular}{@{}p{0.23\columnwidth}p{0.34\columnwidth}p{0.33\columnwidth}@{}}
\toprule
Class & Conditions & Evidence status \\
\midrule
Uniform & 8-bit, 4-bit & Modeled; 4-bit accuracy boundary \\
Selectors & Random, L1, slack-only, signal-only & Measured where local rows exist \\
SUDS main & SUDS+L1, SUDS+signal & Measured accuracy plus modeled PPA \\
SUDS ablation & SUDS-only & Budget-interface isolation \\
Architecture & Lightening, HyAtten & Matched or same-fabric context \\
Boundary fabric & TeMPO, ASTRA & Alternate readout/fabric context \\
\bottomrule
\end{tabular}
\end{table}

The ablation table separates budget allocation from exact-column selection. A
case where signal-only, L1-only, or HyAtten-style selection beats a
SUDS-composed row is kept as boundary evidence that local selectors can
dominate the scheduler-derived budget for that workload and parameter setting.

\section{Calibration Evidence}

The architecture simulator consumes calibration and boundary artifacts rather
than claiming circuit closure. ADC tier energy is tied to the local ADC macro
sanity suite. Sideband-control cost is tied to the Yosys RTL control-plane
artifact. Optical-link assumptions are tied to the PHY boundary sweep. These
artifacts calibrate simulator parameters and stress settings only.

\begin{table*}[t]
\centering
\caption{Parameter traceability for the selected operating point. Local circuit
artifacts are calibration or proxy evidence only. Literature paths are mirrored
in the private knowledge base during development and sanitized in the public
reproduction package.}
\label{tab:calibration}
\footnotesize
\begin{tabular}{llll}
\toprule
Parameter & Value & Evidence label & Claim boundary \\
\midrule
\texttt{tile\_dim} & 32 & literature-anchored assumption & architecture parameter \\
\texttt{tiles} & 4 & literature-anchored assumption & matched reference setting \\
\texttt{cores\_per\_tile} & 2 & literature-anchored assumption & matched reference setting \\
\texttt{frequency\_ghz} & 5.0 & literature-anchored assumption & modeled frequency \\
\texttt{adc8\_pj} & 1.0556 pJ & ADC macro calibration & macro calibration only \\
\texttt{adc6\_pj} & 0.2639 pJ & ADC macro calibration & macro calibration only \\
\texttt{adc4\_pj} & 0.0660 pJ & ADC macro calibration & macro calibration only \\
\texttt{dac\_pj} & 0.6861 pJ & literature-scaled assumption & conversion accounting \\
\texttt{control\_pj\_per\_sideband\_group} & 0.5904 pJ & RTL synthesis proxy & no P\&R or timing closure \\
\texttt{phy\_nominal\_pass\_ratio} & 0.6007 & parametric boundary & optical-link boundary only \\
\texttt{phy\_pessimistic\_laser\_multiplier} & 1.15x & parametric boundary & sensitivity setting \\
\texttt{selected\_sideband\_group\_cols} & 32 & selected design point & RTL anchor \\
\texttt{selected\_adc\_sharing} & temporal accumulation & selected design point & DPTC readout model \\
\bottomrule
\end{tabular}
\end{table*}

\section{Related-Work Boundaries}

The 2026 delta audit did not reopen the selected architecture point, but it did
sharpen the boundary language. ASTRA-style stochastic photonic Transformer
rows are kept as alternate-fabric boundaries. TeMPO-style time multiplexing is
modeled as a boundary readout fabric, not a same-scope SUDS DPTC baseline.
Light-Bound Transformer robustness work is adjacent to noise training rather
than scheduler-derived budget control. Electro-optic attention nonlinearities
target Softmax/Sigmoid-like functions, while this simulator focuses on GEMM,
conversion, memory, optical-link, and control accounting. PRISM-style block
selection targets long-context KV-cache retrieval, and non-photonic CIM
Transformer accelerators remain reviewer-pressure context outside the matched
photonic/DPTC surface.

\section{Reproducibility and Gate Status}

All promoted rows are regenerated through deterministic CSV and JSON artifacts.
The primary command is:
\begin{verbatim}
make suds-optical-transformer-pivot-gate
\end{verbatim}
which rebuilds the architecture simulator, G1 release artifacts, and pivot gate.
The public reproduction manifest includes the TETC architecture simulator
summary, kernel, sensitivity, parameter, design-space, GLUE-linkage, gate,
red-team, and public-repro alignment artifacts. A separate generated public
package check verifies that datasets, weights, private literature mirrors,
trial histories, and personal paths are excluded.

\begin{table}[t]
\centering
\caption{TETC promotion gate after G1 integration. G6 is advisory because
venue partition must be rechecked in the submitter's institutional database.}
\label{tab:gate}
\footnotesize
\begin{tabular}{@{}p{0.31\columnwidth}p{0.13\columnwidth}p{0.44\columnwidth}@{}}
\toprule
Gate & Status & Evidence \\
\midrule
G0 route lock & pass & Protected fallback plus reframe \\
G1 manuscript source & pass & Integrated source, red-team, repro alignment \\
G2 workloads & pass & MobileViT-S and BERT/GLUE on MPS \\
G3 PPA & pass & System-level architecture simulator \\
G4 baselines & pass & Uniform, selector, SUDS, architecture rows \\
G5 calibration & pass & ADC, RTL, PHY tied to parameters \\
G6 venue fit & pass & Advisory route fit record \\
\bottomrule
\end{tabular}
\end{table}

\section{Internal Red-Team Summary}

The G1 red-team artifact uses four lenses: architecture, photonic/circuit,
systems/reproducibility, and reviewer skepticism. The promoted fixes are:
keep SUDS as a scheduler-derived budget interface rather than a selector
replacement; keep ASTRA and TeMPO as boundary fabrics; expose calibration
parameters in Table~\ref{tab:calibration}; keep signal-only and L1-only wins
as boundary evidence; and require public-repro alignment before the gate can
promote. No external reviewer was available inside this local execution pass,
so the internal review is recorded as a substitute, not an equivalent external
review.

\section{Limitations}

This paper does not claim a fabricated photonic chip, extracted layout,
foundry closure, bench-energy measurement, optical-device signoff, or deployed
system. Uniform 4-bit rows remain accuracy boundaries where the accuracy run is
not measured. TeMPO and ASTRA rows are boundary architecture contexts because
they change the conversion or stochastic optical fabric. Full external red-team
review remains preferred before submission, even though the local G1 gate now
has an internal substitute and a public-repro alignment check.

\section{Conclusion}

SUDS reframes scheduler slack as an explicit budget interface for dynamic
photonic Transformer accelerators. When mapped to a Lightening-style DPTC
fabric, the interface composes with L1 and signal selectors, preserves
governed MPS accuracy linkage on BERT-base GLUE and MobileViT-S Transformer
blocks, and maintains a pessimistic system-level EDP advantage versus the
same-scope DPTC reference. The result is a bounded TETC architecture package:
strong enough to promote as architecture evidence, and explicit about where
the evidence stops.

\end{document}
